\chapter*{Введение}
\addcontentsline{toc}{chapter}{Введение}

\section*{Актуальность темы исследования}

% Опора: введение кандидатской (Study/Главы/Автореферат3.docx), обновлённое на
% 2026 г. Ключевой сдвиг относительно 2011 г.: многомерные данные перестали быть
% принадлежностью OLAP-систем и стали повседневной нагрузкой универсальных СУБД
% (геоданные, диапазоны, полнотекстовые и векторные представления), а обобщённое
% дерево поиска стало штатным механизмом расширения СУБД.

\section*{Степень разработанности темы}

% Хеллерстайн (обобщённое дерево поиска), Корнакер (конкурентность и
% восстановление), Гуттман и Бекман (R-дерево, R*-дерево), Бартунов и Сигаев
% (реализация в PostgreSQL). Показать, что нерешённым остаётся жизненный цикл
% индекса: построение, сопровождение, верификация — а не только поиск.

\section*{Цели и задачи}

% Из plan/00-plan.md.

\section*{Научная новизна}

\section*{Теоретическая и практическая значимость}

% Главное отличие работы: все алгоритмы приняты в основную ветку PostgreSQL и
% поставляются с релизами 10--19; эксперименты воспроизводимы на публичных
% версиях, каждый результат идентифицируется коммитом.

\section*{Методология и методы исследования}

\section*{Положения, выносимые на защиту}

\section*{Степень достоверности и апробация результатов}

\section*{Публикации}
